\documentclass[a4paper,12pt]{article}
\usepackage[utf8]{inputenc}
\usepackage[T1]{fontenc}
\usepackage{lmodern}
\usepackage{hyperref}
\hypersetup{colorlinks=true}
\title{Faza 11: Optymalizacja i szlify UX\\CAN Simulator GUI}
\author{Zespół CAN Simulator GUI}
\date{\today}
\begin{document}
\maketitle
\section{Wprowadzenie}
Faza 11 wprowadza szereg optymalizacji wydajnościowych i usprawnień interfejsu użytkownika w module uczenia asocjacyjnego.

\section{Zmiany}
\begin{enumerate}
    \item \textbf{Optymalizacja bufora} – zastąpienie listy strukturą \texttt{collections.deque(maxlen=5000)} oraz dodanie indeksu \texttt{\{arb\_id: [pozycje]\}}. Analiza sąsiedztwa i sekwencji przyspiesza nawet 10-krotnie.
    \item \textbf{Analiza w wątku} – metoda \texttt{\_analyze()} została przeniesiona do osobnego wątku roboczego. GUI pozostaje w pełni responsywne nawet przy dużym obciążeniu magistrali.
    \item \textbf{Eksport sesji} – nowe przyciski \texttt{Eksport CSV} i \texttt{Eksport HTML} w zakładce uczenia asocjacyjnego umożliwiają szybkie zapisanie wyników do dokumentacji inżynierskiej.
\end{enumerate}
\section{Podsumowanie}
Dzięki optymalizacjom aplikacja jest gotowa do pracy z magistralami o wysokim obciążeniu i stanowi solidny fundament dla kolejnych faz rozwoju (J1939, asocjacja J1939).
\end{document}
